%==============================================================================
% Section 2 -- Mathematical Representation of 3D Gaussians
%==============================================================================
\section{Mathematical Representation of 3D Gaussians}
\label{sec:gaussian_repr}

\subsection{The 3D Gaussian Kernel}

A 3D Gaussian is defined by its unnormalised influence kernel:
\begin{equation}
  \boxed{
    \Gcal(\mathbf{x}\,;\,\mean,\cov)
    = \exp\!\left(-\frac{1}{2}(\mathbf{x}-\mean)\transp \cov^{-1} (\mathbf{x}-\mean)\right)
  }
  \label{eq:gaussian_kernel}
\end{equation}
where $\mean \in \R^3$ is the centre and $\cov \in \R^{3\times 3}$ is a
symmetric positive semi-definite (SPSD) covariance matrix. The kernel equals 1
at $\mathbf{x}=\mean$ and decays with the Mahalanobis distance; the locus
$\Gcal = e^{-1/2}$ traces the \emph{one-sigma ellipsoid}.

%------------------------------------------------------------------------------
\subsection{Gaussian Primitive}

\begin{definitionbox}{Gaussian Primitive}
Each primitive $\Gcal_i$ is an oriented, anisotropic ellipsoid characterised by:
\begin{itemize}
  \item $\mean_i \in \R^3$ -- \textbf{position},
  \item $\cov_i \in \R^{3\times 3}$ -- \textbf{shape and orientation}
        (eigenvalues = squared semi-axis lengths; eigenvectors = principal axes),
  \item $o_i \in [0,1]$ -- \textbf{opacity},
  \item $\mathbf{f}_i \in \R^{(L+1)^2 \times 3}$ -- \textbf{colour} as SH
        coefficients (Section~\ref{sec:sh}).
\end{itemize}
\end{definitionbox}

%------------------------------------------------------------------------------
\subsection{Why Gaussians?}

\begin{keybox}{Desirable Properties}
\begin{enumerate}
  \item \textbf{Differentiability}: infinitely differentiable w.r.t.\ all
        parameters, enabling gradient-based optimisation.
  \item \textbf{Closed-form projection}: a 3D Gaussian projects to a 2D
        Gaussian under a perspective camera (Section~\ref{sec:projection}).
  \item \textbf{Soft boundaries}: smooth fall-off models translucent or fuzzy
        regions (hair, smoke) without hard silhouettes.
  \item \textbf{Efficient rasterisation}: 2D splats are rendered in a single
        GPU forward pass via a tile-based rasteriser.
\end{enumerate}
\end{keybox}

%------------------------------------------------------------------------------
\subsection{Scene Representation and Parameter Count}

A scene is a finite mixture of $N$ primitives:
\begin{equation}
  \mathcal{S} = \bigl\{ (\mean_i,\, \cov_i,\, o_i,\, \mathbf{f}_i) \bigr\}_{i=1}^{N}
\end{equation}

Typical scenes contain $N \in [10^5, 10^7]$ Gaussians stored as flat GPU
tensors. For SH degree $L=3$ each Gaussian has 59 learnable parameters
($3 + 4 + 3 + 1 + 48$), so $3\times10^6$ Gaussians require $\approx700$ MB in
\texttt{float32}.

\begin{notebox}
  Directly optimising the 9 entries of $\cov$ does not guarantee SPSD. 3DGS
  instead parametrises $\cov$ through a factorisation that enforces this
  constraint by construction (Section~\ref{sec:covariance}).
\end{notebox}
